%============================================================
% BAB III — METODE PENELITIAN
% Struktur: A. Tempat dan Waktu Penelitian (Bila Ada),
% B. Metode Pengumpulan Data, C. Alat dan Bahan Penelitian,
% D. Konsep Penelitian.
%============================================================
\chapter{METODE PENELITIAN}

\section{Tempat dan Waktu Penelitian (Bila Ada)}
% Ganti teks panduan di bawah ini dengan metode Anda.
Skripsi berupa penelitian mandiri merupakan kajian ilmu
pengetahuan atau teknologi di bidang Teknik Informatika.
Penelitian dapat dilaksanakan dengan mempertimbangkan tingkat
kedalaman dan manfaat penelitian tersebut, khususnya untuk
pengayaan ilmu pengetahuan atau penerapan ilmu pengetahuan dan
teknologi.

\section{Metode Pengumpulan Data}
Bagian ini memuat penjelasan secara lengkap dan terinci tentang
cara-cara yang digunakan dalam proses pengumpulan data untuk jenis
data yang diperlukan. Misalnya melalui observasi, wawancara,
eksperimen, atau kuesioner. Jika metode kuesioner digunakan, maka
blangko angket kuesioner harus dilampirkan dalam laporan.

\section{Alat dan Bahan Penelitian}
Berisi sumber data, software dan hardware yang digunakan untuk
melakukan penelitian skripsi.

\section{Konsep Penelitian}
Konsep penelitian ini digunakan sebagai dasar untuk menjawab
pertanyaan-pertanyaan penelitian yang diangkat. Dapat diartikan
sebagai mengalirkan jalan pikiran menurut kerangka logis
(\emph{construct logic}) atau kerangka konseptual yang relevan
untuk menjawab penyebab terjadinya masalah (tidak harus membuat
sistem). Jika membuat produk maka menggunakan metode pengembangan
sistem, jika tidak maka alur disesuaikan dengan penelitian yang
akan dilakukan, misalnya menggunakan metode SPK, algoritma dalam
data mining, framework tatakelola sistem informasi, dan lain-lain.
